% Part: many-valued-logic
% Chapter: syntax-and-semantics
% Section: matrices

\documentclass[../../../include/open-logic-section]{subfiles}

\begin{document}

\olfileid[hi]{mvl}{syn}{mat}

\olsection{आव्यूह}

कोई बहुमूल्य तर्कशास्त्र अपनी भाषा, सत्य-मानों के समुच्चय~$V$, निर्दिष्ट
सत्य-मानों के उपसमुच्चय और अपने संयोजकों के सत्य-फलनों से परिभाषित होता
है। इन अवयवों को एक साथ \emph{आव्यूह} कहते हैं।

\begin{defn}[आव्यूह]
\ollabel{defn:matrix}
तर्कशास्त्र~$\Log L$ के किसी \emph{आव्यूह} में ये होते हैं:
\begin{enumerate}
\item भाषा~$\Lang L$ बनाने वाला संयोजकों का समुच्चय;
\item सत्य-मानों का समुच्चय $V \neq \emptyset$;
\item निर्दिष्ट सत्य-मानों का समुच्चय $V^+ \subseteq V$;
\item ~ $\Lang L$ के प्रत्येक $n$-स्थानी संयोजक $\star$ के लिए सत्य-फलन
~$\tf{\star} : V^n \to V$। यदि $n = 0$, तो $\tf{\star}$ केवल~$V$ का कोई
अवयव है।
\end{enumerate}
\end{defn}

\begin{ex}
चिरसम्मत तर्कशास्त्र~\LogCL{} के आव्यूह में ये होते हैं:
\begin{enumerate}
  \item मानक प्रतिज्ञप्तिपरक भाषा $\Lang L_0$, जिसमें
  $\lfalse$, $\lnot$, $\land$, $\lor$, $\lif$.
  \item सत्य-मानों का समुच्चय $V = \{\True, \False\}$।
  \item $\True$ एकमात्र निर्दिष्ट मान है, अर्थात् $V^+ = \{\True\}$।
  \item $\lfalse$ के लिए $\tf{\lfalse} = \False$। अन्य सत्य-फलन प्रचलित
  सत्य-सारणियों से दिए जाते हैं (\olref{fig:tf-CL} देखें)।
\end{enumerate}
\begin{figure}
  \begin{center}
      \begin{tabular}{c|c} 
        $\tf{\lnot}$ & \\ 
        \hline  
        $\True$ & $\False$ \\ 
        $\False$ & $\True$ 
      \end{tabular}
      \quad
      \begin{tabular}{c|cc} 
        $\tf{\land}$ & $\True$ & $\False$ \\ 
        \hline 
        $\True$ & $\True$ & $\False$ \\ 
        $\False$ & $\False$ & $\False$ 
      \end{tabular}
      \quad
      \begin{tabular}{c|cc} 
        $\tf{\lor}$ & $\True$ & $\False$ \\ 
        \hline 
        $\True$ & $\True$ & $\True$ \\ 
        $\False$ & $\True$ & $\False$ 
      \end{tabular}
      \quad
      \begin{tabular}{c|cc} 
        $\tf{\lif}$ & $\True$ & $\False$ \\ 
        \hline 
        $\True$ & $\True$ & $\False$ \\ 
        $\False$ & $\True$ & $\True$ 
      \end{tabular}
    \end{center} 
    \caption{चिरसम्मत तर्कशास्त्र~$\LogCL$ के सत्य-फलन।}
    \ollabel{fig:tf-CL}
  \end{figure}
  \end{ex}

\end{document}
